% Dieselbe Definitionsformel wie in Band 28 anzeigen, ohne neue Registrierung.
\newcommand{\PowerProductDefinitionFormula}{%
  \begin{aligned}[t]
    &\SemigroupAlg{A}{\star}\dsep X,Y\in\PowNE{A}\vdash{}\\[-2pt]
    &X\star_{\mathcal P}Y\coloneqq
      \{x\star y\mid x\in X,\ y\in Y\}
  \end{aligned}%
}

% Die Originalaussagen anzeigen, ohne sie in der Ergänzung erneut zu registrieren.
\newcommand{\PowerProofStatement}[1]{%
  \par\Needspace{15\baselineskip}%
  \hypertarget{power.proof.#1}{}%
  \begingroup
  \RenewDocumentCommand{\FormulaThmDeltaK}{o m m +m}{%
    \noindent\textbf{Beweis zu \FormulaRefAuto{##3}[thm]}%
    \IfNoValueF{##1}{\space(##1)}\par
    {\small\begin{DeltaContext}{ }##4\end{DeltaContext}}%
    \[##2\]%
  }%
  \PowerStatement{#1}%
  \endgroup
}

% Eine bereits in Band 28 registrierte Hilfsaussage mit eigenem Beweis zeigen.
\newcommand{\PowerInnerProofStatement}[1]{%
  \par\Needspace{15\baselineskip}%
  \hypertarget{power.proof.#1}{}%
  \noindent\textbf{Beweis zu \FormulaRefAuto{#1}[thm]}%
  \space(\csname PowerInnerTitle@#1\endcsname\unskip)\par
  \[\csname PowerInnerFormula@#1\endcsname\]%
}

% Ergebniszuordnungen bleiben als Quellmarkierungen erhalten; die Tabellen
% zeigen ausschließlich die mathematischen Aussagen.
\NewDocumentCommand{\PowerInnerProofResult}{O{} m +m}{#3}
